\par\bigskip\Needspace{9\baselineskip}
\pdfbookmark[2]{Zweite Methode, [6], Olbers, 2. Julius 1805}{ceres-second-6}
\begin{center}[6.]\par\medskip
\textsc{Gauss} an \textsc{Olbers}, Braunschweig, 2. Julius 1805.
\end{center}
\noindent\hspace*{\parindent}. . . . .\quad Mit meiner Rechnung der Störungen der {\AstroSymbols ⚳} bin ich leider noch nicht fertig. Theils habe ich nicht immer anhaltend daran gearbeitet, theils habe ich auch daran bei weitem mehr Mühe gehabt, als ich vorher glaubte, hie und da auch wohl mehr, als nothwendig gewesen wäre. Aber die zweckmässigste Ausübung einer Methode lernt man erst bei ihrer Anwendung. Mit den Störungen der Breite hatte ich angefangen. Ich hatte bei meiner ersten Rechnung alles, was von den Excentricitäten abhing, übergangen: \textsc{Oriani} hatte nur Ein Glied mitgenommen. Nunmehro habe ich alle Gleichungen, die ich über \(1''\) fand, mitgenommen, worunter auch ein paar sind, die von dem Producte der Excentricitäten abhangen, also (da ohnehin in alle Breitengleichungen die Neigung entrirt) von der Ordnung 3 sind. Von dieser Arbeit habe ich nun gewissermassen auch schon Früchte geerndtet. Sie erinnern sich, dass ich mich schon seit 1803 beklagt habe, dass sich die Breiten in der {\AstroSymbols ☍} von 1803 mit denen von 1801 nach der aus der Theorie gefundenen Bewegung des \(\AscendingNode\) nicht mehr vereinigen liessen, und dass ich gezwungen war, den \(\AscendingNode\) 1803 um \(3'\) weiter zu rücken. Eben so stimmte in der {\AstroSymbols ☍} 1804 die Breite nicht mit der von 1802; die Neigung musste weit mehr verringert werden, als die Theorie angab.

Bei meinen neuesten Elementen musste ich dem \(\AscendingNode\) eine tägliche tropische Bewegung von \(0''{,}241\) und der Neigung eine tägliche Abnahme von \(0''{,}0243\) geben, um die Beobb. zu vereinigen, so dass
\begin{center}\small
\begin{tabular}{@{}l|l@{}}
\(\AscendingNode\) 1801 Jan. 1. \(80^{0}\,54'\,46''\)&Neigung 1802 {\AstroSymbols ☍} \(10^{0}\,38'\ \ 1''\)\\
\hspace*{1em}1803\quad{\AstroSymbols ☍}\quad\(80\ \ 58\ \ 28\)&\hspace*{4em}1804 {\AstroSymbols ☍} \(10\ \ 37\ \ 38.\)
\end{tabular}
\end{center}

\SourceLeaf{414}Zu meiner grossen Freude ist dies nun nicht mehr nöthig, und die Beobachtungen stimmen jetzt mit den neuen Breitengleichungen und Säcularbewegungen recht gut. Die an sich unbedeutenden noch zurückbleibenden Differenzen von ein paar Secunden können recht gut dadurch erklärt werden, dass die dabei gebrauchten Radii vectores noch Verbesserung nöthig haben, weil dabei die Störungsgleichungen noch nicht vollständig angewandt sind. Hier die Breitenstörungen:
\begin{center}\small
\begin{tabular}{@{}lllll@{}}
{\AstroSymbols ⚳}:&Mittl.&Länge der {\AstroSymbols ⚳},&\(\tilde\omega\):&Sonnenferne\\
{\AstroSymbols ♃}:&»&»\quad des {\AstroSymbols ♃},&\(\tilde\omega'\):&»
\end{tabular}\par\smallskip
\(\AscendingNode\) Länge des aufst. Knotens der {\AstroSymbols ⚳} auf der {\AstroSymbols ♃}-Bahn \(=78^{0}\,35'\)
\end{center}
\begin{center}\small
\begin{tabular}{@{}llr@{}}
Knoten der {\AstroSymbols ⚳} 1801&\(.\ .\ .\ .\ .\ .\ .\ .\ .\ .\ .\ .\)&\(80^{0}\,53'\,56''\)\\
Tägl. Bewegung&\(.\ .\ .\ .\ .\ .\ .\ .\ .\ .\ .\ .\ .\ .\ .\ .\)&\(+0''{,}00405\)\\
Neigung 1801&\(.\ .\ .\ .\ .\ .\ .\ .\ .\ .\ .\ .\ .\ .\ .\ .\ .\ .\)&\(10^{0}\,37'\,29''{,}9\)\\
Tägl. Abnahme&\(.\ .\ .\ .\ .\ .\ .\ .\ .\ .\ .\ .\ .\ .\ .\ .\ .\ .\)&\(0''{,}001195.\)
\end{tabular}
\end{center}
\begin{center}Breitengleichungen:\par\medskip A.\end{center}
\begingroup\small
\[
\renewcommand{\arraystretch}{1.18}
\begin{array}{@{}l@{\qquad}l@{}}
\left.\begin{aligned}&-11''{,}19\sin(\text{\AstroSymbols ♃}-\AscendingNode)\\&\underline{+\phantom{1}2{,}78\sin(\tilde\omega-\AscendingNode)\cos(\text{\AstroSymbols ♃}-\tilde\omega')}\end{aligned}\right\}&-9''{,}65\sin(\text{\AstroSymbols ♃}-74^{0}\,13')\\[5pt]
+14{,}17\sin(\text{\AstroSymbols ⚳}-2\text{\AstroSymbols ♃}+\AscendingNode)&+14{,}17\sin(\text{\AstroSymbols ⚳}-2\text{\AstroSymbols ♃}+78^{0}\,35')\\
+27{,}47\sin(2\text{\AstroSymbols ⚳}-3\text{\AstroSymbols ♃}+\AscendingNode)&+27{,}47\sin(2\text{\AstroSymbols ⚳}-3\text{\AstroSymbols ♃}+78^{0}\,35')\\
-\phantom{1}4{,}62\sin(3\text{\AstroSymbols ⚳}-4\text{\AstroSymbols ♃}+\AscendingNode)&-\phantom{1}4{,}62\sin(3\text{\AstroSymbols ⚳}-4\text{\AstroSymbols ♃}+78^{0}\,35')\\
-\phantom{1}0{,}01\sin(4\text{\AstroSymbols ⚳}-5\text{\AstroSymbols ♃}+\AscendingNode)&-\phantom{1}0{,}01\sin(4\text{\AstroSymbols ⚳}-5\text{\AstroSymbols ♃}+78^{0}\,35')\\
+\phantom{1}5{,}51\sin(2\text{\AstroSymbols ⚳}-\text{\AstroSymbols ♃}-\AscendingNode)&+\phantom{1}5{,}51\sin(2\text{\AstroSymbols ⚳}-\text{\AstroSymbols ♃}-78^{0}\,35')\\
+\phantom{1}0{,}99\sin(3\text{\AstroSymbols ⚳}-2\text{\AstroSymbols ♃}-\AscendingNode)&+\phantom{1}0{,}99\sin(3\text{\AstroSymbols ⚳}-2\text{\AstroSymbols ♃}-78^{0}\,35')
\end{array}
\]
\endgroup
\clearpage
\begin{center}B.\end{center}
\begingroup\small
\[
\renewcommand{\arraystretch}{1.10}
\begin{array}{@{}l@{\qquad}l@{}}
\underline{-2''{,}71\sin(\tilde\omega-\AscendingNode)}&+1''{,}91\\[3pt]
\left.\begin{aligned}&\underline{+3{,}44\sin(\text{\AstroSymbols ⚳}-\text{\AstroSymbols ♃}-\tilde\omega+\AscendingNode)}\\&\underline{-3{,}33\sin(\text{\AstroSymbols ⚳}-\text{\AstroSymbols ♃}+\tilde\omega-\AscendingNode)}\\&\underline{+2{,}25\sin(\text{\AstroSymbols ⚳}-\text{\AstroSymbols ♃}+\tilde\omega'-\AscendingNode)}\\&\underline{-0{,}24\sin(\text{\AstroSymbols ⚳}-\text{\AstroSymbols ♃}-\tilde\omega'+\AscendingNode)}\end{aligned}\right\}&+8{,}61\sin(\text{\AstroSymbols ⚳}-\text{\AstroSymbols ♃}+95^{0}\,26')\\[3pt]
\SourceLeaf{415}\left.\begin{aligned}&\underline{-4''{,}59\sin(2\text{\AstroSymbols ⚳}-2\text{\AstroSymbols ♃}-\tilde\omega+\AscendingNode)}\\&\underline{+3{,}59\sin(2\text{\AstroSymbols ⚳}-2\text{\AstroSymbols ♃}+\tilde\omega-\AscendingNode)}\\&\underline{-2{,}93\sin(2\text{\AstroSymbols ⚳}-2\text{\AstroSymbols ♃}+\tilde\omega'-\AscendingNode)}\end{aligned}\right\}&-10''{,}31\sin(2\text{\AstroSymbols ⚳}-2\text{\AstroSymbols ♃}+94^{0}\,10')\\[3pt]
\left.\begin{aligned}&\underline{-2{,}26\sin(3\text{\AstroSymbols ⚳}-3\text{\AstroSymbols ♃}-\tilde\omega+\AscendingNode)}\\&\underline{+3{,}90\sin(3\text{\AstroSymbols ⚳}-3\text{\AstroSymbols ♃}+\tilde\omega'-\AscendingNode)}\end{aligned}\right\}&-6{,}16\sin(3\text{\AstroSymbols ⚳}-3\text{\AstroSymbols ♃}+112^{0}\,26')\\[3pt]
\left.\begin{aligned}&\underline{+\phantom{1}4{,}17\sin(2\text{\AstroSymbols ⚳}-\AscendingNode-\tilde\omega')}\\&\underline{-10{,}41\sin(2\text{\AstroSymbols ⚳}-\AscendingNode-\tilde\omega)}\end{aligned}\right\}&-13{,}69\sin(2\text{\AstroSymbols ♃}-57^{0}\,21')\\[3pt]
\underline{-1{,}65\sin(3\text{\AstroSymbols ⚳}-\AscendingNode-\tilde\omega-\tilde\omega')}&-1{,}65\sin(3\text{\AstroSymbols ♃}-56^{0}\,9')\\[3pt]
\left.\begin{aligned}&\underline{+2{,}72\sin(3\text{\AstroSymbols ♃}-\text{\AstroSymbols ⚳}-\AscendingNode-\tilde\omega')}\\&\underline{-4{,}56\sin(3\text{\AstroSymbols ♃}-\text{\AstroSymbols ⚳}-\AscendingNode-\tilde\omega)}\end{aligned}\right\}&-6{,}77\sin(3\text{\AstroSymbols ♃}-\text{\AstroSymbols ⚳}-61^{0}\,25')\\[3pt]
\left.\begin{aligned}&\underline{-2{,}81\sin(2\text{\AstroSymbols ⚳}-4\text{\AstroSymbols ♃}+\AscendingNode+\tilde\omega')}\\&\underline{+1{,}82\sin(2\text{\AstroSymbols ⚳}-4\text{\AstroSymbols ♃}+\AscendingNode+\tilde\omega)}\end{aligned}\right\}&+4{,}30\sin(2\text{\AstroSymbols ⚳}-4\text{\AstroSymbols ♃}+72^{0}\,24')\\[3pt]
\left.\begin{aligned}&+13{,}75\sin(3\text{\AstroSymbols ⚳}-5\text{\AstroSymbols ♃}+\AscendingNode+\tilde\omega')\\&-10{,}56\sin(3\text{\AstroSymbols ⚳}-5\text{\AstroSymbols ♃}+\AscendingNode+\tilde\omega)\end{aligned}\right\}&-22{,}51\sin(3\text{\AstroSymbols ⚳}-5\text{\AstroSymbols ♃}+70^{0}\,27').
\end{array}
\]
\endgroup

Hiemit ist \textsc{Orianis} Rechnung im Dec. 1802 der M. C. zu vergleichen. Alle roth unterstrichenen\HistoricalNoteEight{Die Striche sind im Abdruck wiedergegeben worden.]} Gleichungen sind von ihm übergangen. Der numerische Werth ist
\begin{center}\small
\setlength{\tabcolsep}{6pt}
\begin{tabular}{@{}lrrr@{}}
&\shortstack{von den A bez.\\Gleichg.}&\shortstack{von den\\Gleichg. B}&von allen\\[2pt]
Jan. 1.\quad1801&\(-47''{,}03\)&\(-17''{,}14\)&\(-64''{,}17\)\\
Febr. 11. 1801&\(-45{,}88\)&\(-17{,}56\)&\(-63{,}44\)\\
\hspace*{3em}{\AstroSymbols ☍}\quad1802&\(-25{,}78\)&\(+27{,}20\)&\(+\phantom{2}1{,}41\)\\
\hspace*{3em}{\AstroSymbols ☍}\quad1803&\(+\phantom{2}0{,}46\)&\(+41{,}53\)&\(+41{,}99\)\\
\hspace*{3em}{\AstroSymbols ☍}\quad1804&\(+22{,}87\)&\(-33{,}60\)&\(-10{,}73.\)
\end{tabular}
\end{center}

Hieraus ist klar, dass die Vernachlässigung der Gleichungen B dem alleinigen Gebrauche der A allen Werth nahm.
\begin{center}\rule{.30\linewidth}{.4pt}\end{center}
